\chapter{Aplikacje do programowania na telefonie}

Programowanie to jedna z najpopularniejszych aktywności informatyków naszego pokolenia na telefonie. Każda osoba w naszej grupie przynajmniej raz stała w korku wracając autobusem (na przykład z nie etycznej lekcji etyki) i zastanawiała się co by było gdyby mogła w tym momencie wyjąć telefon i zacząć programować. Innym przypadkiem kiedy aplikacja była by użyteczna jest moment kiedy chcemy stworzyć prosty prototyp strony a jednocześnie nie chcemy wyjść z łóżka, zamiast zabierać komputer do łóżka jak psychopata wygodniej było by sięgnąć po telefon i na szybko coś napisać. Są też ludzie (głównie dzieci) którzy nie mają dostępu do komputera z różnych powodów a i tak chcieli by coś tworzyć.
\newline\newline
Czy naprawdę powinniśmy tym biednym ludziom odbierać możliwość programowania. Naszym zdaniem nie, szczególnie w obecnej erze gdzie każdy może sobie wygenerować szczochy wykorzystując sztuczną inteligencje, powinniśmy wspierać ludzi którzy wciąż mimo przeciwności losu chcą stworzyć coś własnego zamiast iść na łatwiznę. Narządzie dla ludzi z prawdziwą pasją powinny być proste oraz użyteczne a naszą misją jest stworzenie właśnie takiej aplikacji.